\section{Related Work}
\label{sec:related}

\subsection{Python Native-Extension Analysis}

Pungi models Python/C reference-count operations as affine programs
\cite{pungi2014}; CpyChecker tracks objects and expected counts over GCC paths
\cite{cpychecker}; and PyRefcon adds lifecycle states, nullity, and C++ monitors
\cite{pyrefcon2023}. PyCGuard instead covers 13 API obligation categories and
represents correctness as alternative discharge conditions. PyRefcon remains
the closest reproducible baseline because its lifecycle scope overlaps
with the reference ownership, reference validity, and post-release
safety obligations.

Jinn synthesizes dynamic state-machine checkers for JNI and Python/C rules
\cite{jinn2010}. Other work characterizes Python/C API evolution and bugs
\cite{hu2020pythoncapi}, jointly analyzes Python and C with shared abstract
domains \cite{monat2021multilanguage}, or diagnoses inefficient boundary
interactions \cite{pieprof2021}. PyCGuard targets statically visible native-side
API duties rather than executed violations, language interaction errors, or
performance.

General reference-count work uses shallow-alias analysis, compositional
verification, path-pair inconsistency, or two-dimensional consistency
\cite{lal2010refcount,referee2009,rid2016,cid2021}. Resource analyses detect
path-sensitive leaks, infer acquire/release pairs, track C++ smart-pointer heap
protocols, or identify disordered cleanup
\cite{smoke2019,sinkfinder2020,smartpointer2021,hero2021}. These techniques
inform PyCGuard's path and resource reasoning but omit CPython ownership,
exception-state, garbage-collection, and API-specific discharge semantics.

\subsection{Cross-Language Interface Analysis}

JNI analyses detect resource and typestate errors \cite{kondoh2008jni}, analyze
exceptional paths \cite{li2009jni,li2014jniexception}, prune before fine-grained
exception checking \cite{jet2011}, or infer types across FFI boundaries
\cite{furr2008ffi}. They demonstrate the need for native-side models but each
targets a narrower family of properties.

Related systems detect type and lifetime errors in JavaScript bindings
\cite{brown2017jsbindings}, coordinate garbage collection across components
\cite{degenbaev2018crosscomponentgc}, or track polyglot vulnerability flows
\cite{charon2025}. PyCGuard checks CPython API conformance even when both trigger
and consequence remain in native code.

RustBelt formalizes unsafe abstractions \cite{rustbelt2018}; empirical studies
characterize unsafe usage and requirements
\cite{unsaferustoopsla2020,rustsafetypldi2020,unsafeachilles2024}; and Rudra,
MirChecker, SafeDrop, and subsequent analyses detect unsafe-code patterns and
safety bugs \cite{rudra2021,mirchecker2021,safedrop2023,rustsafetytse2024,dontpanic2025}.
CPython lacks an equivalent language-level discipline, so PyCGuard recovers
heterogeneous protocols from documentation and bugs.

\subsection{LLM-Assisted Bug Detection}

Learned vulnerability detectors are sensitive to dataset construction and
generalization \cite{chakraborty2022deepvuln,steenhoek2023vuln}, and standalone
LLMs remain unreliable vulnerability reasoners \cite{ullah2024llmvuln}.
GPTScan instead combines GPT with program analysis \cite{gptscan2024}, supporting
PyCGuard's use of explicit program facts and properties.

Hybrid systems validate static warnings \cite{li2024llmstatic}, infer taint
specifications for whole-repository analysis \cite{iris2025}, couple LLMs with
bidirectional taint analysis \cite{staint2025}, or derive binary-analysis
semantics \cite{latte2025}. PyCGuard starts from an audited Python/C obligation
vocabulary and delegates only candidate path adjudication to the model.

Data-flow facts can constrain LLM bug judgments \cite{wang2024dataflowbug};
LLMDFA decomposes data-flow reasoning and externally validates path feasibility
\cite{llmdfa2024}; DAInfer extracts aliasing specifications from documentation
\cite{dainfer2024}; and KNighter synthesizes static checkers from patches
\cite{knighter2025}. PyCGuard likewise narrows reasoning, but keeps obligations
and pruning rules explicit and auditable.

Safe4U decomposes unsafe Rust contracts before checking wrappers
\cite{safe4u2025}, while RFCAudit and RFCParser interpret protocol documents
\cite{rfcaudit2025,rfcparser2025}. PyCGuard instead normalizes many API statements
and real bugs into reusable obligation types with alternative discharge
conditions and separate pruning patterns.
